\scanpage{13}
\mspage{3}

forme initiale (cas $\alpha,\beta$) implique Hyp. 2 sous sa forme générale (en termes de $\alpha_0,\alpha_1,\ldots\alpha_n$).

\uline{Hyp.} 2': Pour tout $\alpha_0\in\mathrm{Ob}\,A$, le foncteur induit $A_{/\alpha_0}\unsure{\xrightarrow{\varphi}|\twoheadrightarrow|\longrightarrow}\mathcal X_{/\alpha_0}$ satisfait \unsure{à|:} Hyp 2 i.e. pour tout diagramme
\diagram{\begin{tikzcd}[ampersand replacement=\&,column sep=small,row sep=small]
\alpha_1\arrow[dr]\&\&\alpha_2\arrow[dl]\\
\&\alpha_0\&
\end{tikzcd}}
, regardant les carrés commutatifs de $A$ qui le complètent en
\diagram{\begin{tikzcd}[ampersand replacement=\&,column sep=small,row sep=small]
\&\gamma\arrow[dl]\arrow[dr]\&\\
\alpha_1\arrow[dr]\&\&\alpha_2\arrow[dl]\\
\&\alpha_0\&
\end{tikzcd}}
, i.e. les objets de $A_{/\alpha_1\times_{\alpha_0}\alpha_2}$, la famille des flèches correspondantes
\[U_\gamma\longrightarrow U_{\alpha_1}\times_{U_{\alpha_0}}U_{\alpha_2}\]
est couvrante (i.e., épimorphique).

Les hypothèses Hyp 1, 2, 2' se reformulent joliment, en termes du prolongement canonique
\[\varphi_!:\widehat A\longrightarrow\mathcal X\]
de $\varphi$ en un foncteur \inserted{$\widehat A\longrightarrow\mathcal X$} commutant \deleted{\unsure{à|et}} aux petites $\lim_{\longrightarrow}$. Ce sont respectivement
\margin{12}
